\chapter{Crypto High-Frequency Trading on Centralised Venues}\label{ch:hf:crypto-high-frequency-trading-on-centralised-venues}

Bitcoin trades on dozens of exchanges around the clock with no consolidated tape. A market maker quotes on one, hedges on another, pays funding
every eight hours on a third, and has to fit it all inside each venue's request budget. In this chapter's simulated week, a maker quoting a
perpetual at two basis points around another venue's spot price earns \$575\,000 of spread, pays \$464\,000 to hedge and collects \$144\,000 of maker
rebates; a liquidation cascade adds \$20\,800 if it absorbs the forced sells at the right depth. A request budget that allows only one requote every
four seconds costs it \$102\,000 of that week.

\section{Cross-exchange making and taking}

Crypto's venues (One Quant Book 3, chapters 15 and 16) are separate order books with separate custody, fees and rules; the same coin trades on all
of them, and nothing routes an order from one to the best price on another. Makarov and Schoar found large and recurrent arbitrage opportunities
across exchanges, much larger across countries than within them, and attributed their persistence to capital controls and the difficulty of moving
funds; within a country, prices stay close because traders like the ones this chapter describes keep them there.

\begin{definition}[Cross-exchange market making]\label{def:hf:crypto-high-frequency-trading-on-centralised-venues:xmm}
\emph{Cross-exchange market making}\index{cross-exchange market making} is quoting on one venue at prices derived from another venue's book, and
hedging each fill at once by taking liquidity on the other venue, so that the maker earns the difference between the spread it quotes and the cost of
the hedge.
\end{definition}

\begin{definition}[Hedge venue]\label{def:hf:crypto-high-frequency-trading-on-centralised-venues:hedge}
A \emph{hedge venue}\index{hedge venue} is the exchange on which a cross-exchange market maker offsets the positions its quotes accumulate on
another: usually the deepest book for the coin, where its takes move the price least and whose price the maker uses as its fair value.
\end{definition}

\begin{omfigure}
\begin{tikzpicture}[font=\footnotesize,>=stealth,box/.style={draw,rounded corners,fill=omThm!8,minimum width=2.6cm,minimum height=0.8cm,align=center}]
\node[box] (a) at (0,0) {quote venue\\perpetual};
\node[box,fill=omDef!10] (m) at (4.5,0) {market maker\\(fair = spot)};
\node[box] (b) at (9,0) {hedge venue\\spot book};
\node[box,fill=gray!10] (c) at (0,-2.1) {clients, arbitrageurs,\\forced sells};
\node[box,fill=gray!10] (fu) at (4.5,-2.1) {funding\\every 8 hours};
\draw[<->] (m) -- node[above] {quotes, 2 bp} (a);
\draw[->] (b) -- node[above] {fair price} (m);
\draw[->] (m.south east) to[bend right=15] node[below right] {hedge: take} (b.south west);
\draw[->] (c) -- node[left] {fills} (a);
\draw[->] (fu) -- (m);
\node[font=\scriptsize,align=center] at (4.5,1.2) {request budget: one requote\\every $1/r$ seconds};
\end{tikzpicture}

\omcaption{A cross-exchange market maker: it quotes the perpetual on one venue around the \omterm{def:hf:fair-value:fair}{fair price} read from the spot book of another, hedges
each fill there by taking, receives or pays funding on its perpetual position every eight hours, and can move its quotes only as often as its
request budget allows. Source: this chapter.}\label{fig:hf:crypto-high-frequency-trading-on-centralised-venues:flow}
\end{omfigure}

The chapter's week (\cref{fig:hf:crypto-high-frequency-trading-on-centralised-venues:flow}): a coin at \$60\,000 with 3\% daily volatility; clients
send \$10\,000 orders to the perpetual half a second apart on average, 52\% of them buys; arbitrageurs take \$20\,000 whenever a quote has gone stale
by more than its half-spread plus a basis point. The maker quotes 2 basis points either side of the \omterm{def:hf:fair-value:fair}{fair price}, pays a maker fee of $-0.5$ basis
point (a rebate), and hedges on spot by taking at a half-spread of 0.5 basis point plus a 1-basis-point fee and 20 milliseconds of slippage. It holds
at most \$200\,000 of perpetual inventory on either side (the hedge makes it delta-neutral; the limit bounds the basis risk).

\section{Funding and basis at speed}

A perpetual future (One Quant Book 3, chapter 17) has no expiry; a funding payment at fixed times keeps it close to the index. A maker whose
perpetual position is hedged in spot is exposed not to the coin but to the basis between the two and to the funding paid on the perpetual leg. In the
model the clients' slight preference for buying leaves the maker short the perpetual (long spot) most of the week; with funding at 1 basis point
per interval on average, it receives a few hundred dollars. Funding matters more to the maker's timing than to its level: positions carried across
a funding timestamp pay or earn the interval's rate, positions opened after it do not, and the rate is known in advance from the premium index.

\section{Liquidation flows}

\begin{definition}[Liquidation absorption]\label{def:hf:crypto-high-frequency-trading-on-centralised-venues:absorb}
\emph{Liquidation absorption}\index{liquidation absorption} is providing liquidity to the forced orders a venue's liquidation engine sends into its
book during a cascade, at prices that discount the forced selling, with a hedge on another venue and the capacity to hold the position until the
dislocation reverts.
\end{definition}

The planted cascade: on the fourth day at noon the perpetual falls 80 basis points below spot over five minutes and recovers over the next hour, with a ten-minute time constant
(\cref{fig:hf:crypto-high-frequency-trading-on-centralised-venues:cascade}); forced sells of \$50\,000 a second arrive for five minutes. A maker that
pulls its bid while the forced sells arrive earns nothing from it (\$600 over the week's cascade account). One that absorbs them into a \$5 million
limit from the first second buys early, at shallow discounts, and earns \$5\,200 from the basis's recovery; one that waits until the perpetual is
40 basis points below spot earns \$20\,800; at 60 basis points it is too late to fill the limit, and earns \$17\,800.

\begin{omfigure}
\begin{tikzpicture}
\begin{axis}[width=11cm,height=5.4cm,axis y line*=right,axis x line=none,xmin=-2,xmax=60,ymin=-90,ymax=10,
  ylabel={\small perpetual less spot (bp)},tick label style={font=\footnotesize},table/col sep=comma]
\addplot[gray,thick] table[x=minute,y=basis]{figdata/hft/24-crypto-high-frequency-trading-on-centralised-venues/cascade.csv};
\label{plot:hf:cbasis}
\end{axis}
\begin{axis}[width=11cm,height=5.4cm,axis y line*=left,xmin=-2,xmax=60,xlabel={\small minutes from the start of the cascade},
  ylabel={\small perpetual position (\$ million)},ymin=-0.5,ymax=5.5,tick label style={font=\footnotesize},grid=major,grid style={gray!25},
  table/col sep=comma,legend style={font=\scriptsize,at={(0.5,-0.3)},anchor=north},legend columns=3]
\addplot[omThm,thick] table[x=minute,y=absorb]{figdata/hft/24-crypto-high-frequency-trading-on-centralised-venues/cascade.csv};
\addlegendentry{absorbing from $-40$ bp}
\addplot[omDef,thick,dashed] table[x=minute,y=pull]{figdata/hft/24-crypto-high-frequency-trading-on-centralised-venues/cascade.csv};
\addlegendentry{pulling the bid}
\addlegendimage{/pgfplots/refstyle=plot:hf:cbasis}\addlegendentry{basis (right)}
\end{axis}
\end{tikzpicture}

\omcaption{The planted liquidation cascade: the perpetual's basis to spot (right axis) and the maker's perpetual position, hedged in spot (left
axis), for a maker that absorbs the forced sells once the basis is 40 basis points deep, up to \$5 million, and for one that pulls its bid. Data:
\texttt{hf\_crypto.absorption}.}\label{fig:hf:crypto-high-frequency-trading-on-centralised-venues:cascade}
\end{omfigure}

\begin{dated}{2026-09}{Limits, funding clocks and the October 2025 cascade}\label{dat:hf:crypto-high-frequency-trading-on-centralised-venues:limits}
A large spot venue's exchange-information endpoint returned, on 24 September 2026, request-weight limits of 6\,000 a minute and order limits of 100 per
10 seconds and 200\,000 a day. The same venue's futures funding is exchanged at 00:00, 08:00 and 16:00 UTC, with an interest component of 0.01\% per
interval for most contracts. On 10 October 2025, after a tariff announcement, \$19 billion of leveraged crypto positions were liquidated within 24
hours, as a research report of a crypto news publisher counted them (One Quant Book 3, chapter 18).
\end{dated}

\section{Rate-limit-aware quoting}

A crypto venue meters requests (One Quant Book 3, chapter 15): weights per minute, orders per ten seconds, orders per day. The tightest rule sets
how often the maker can move its quotes. With the dated box's limits and one instrument, the daily order cap allows 2.3 orders a second, 1.16
requotes of two orders each: the daily cap, not the ten-second one, binds. Quoting five instruments within the same budget leaves each 0.23 requotes a
second.

\begin{omfigure}
\begin{tikzpicture}
\begin{axis}[width=11cm,height=5.2cm,xlabel={\small requotes a second},ylabel={\small \$ thousand over the week},xmode=log,log basis x=10,
  xtick={0.1,0.2,0.25,0.5,1},xticklabels={0.1,0.2,0.25,0.5,1},xmin=0.09,xmax=1.1,ymin=-700,ymax=320,tick label style={font=\footnotesize},
  grid=major,grid style={gray!25},table/col sep=comma,legend style={font=\scriptsize,at={(0.97,0.05)},anchor=south east}]
\addplot[omThm,thick,mark=*,mark size=1.4pt] table[x=rate,y=total]{figdata/hft/24-crypto-high-frequency-trading-on-centralised-venues/budget.csv};
\addplot[omDef,thick,dashed,mark=square*,mark size=1.3pt] table[x=rate,y=adverse]{figdata/hft/24-crypto-high-frequency-trading-on-centralised-venues/budget.csv};
\draw[gray,dashed] (axis cs:0.09,0) -- (axis cs:1.1,0);
\legend{week's total,picked off by arbitrageurs}
\end{axis}
\end{tikzpicture}

\omcaption{The maker's total over the simulated week and the part lost to arbitrageurs taking stale quotes, against how often its request budget
lets it move its quotes (log scale): once a second, every two, four, five and ten seconds. Data:
\texttt{hf\_crypto.budget}.}\label{fig:hf:crypto-high-frequency-trading-on-centralised-venues:budget}
\end{omfigure}

The cost of a thin budget is stale quotes (\cref{fig:hf:crypto-high-frequency-trading-on-centralised-venues:budget}): at one requote a second the
arbitrageurs find nothing; every two seconds they take \$6\,400 over the week; every four, \$135\,000; every ten, \$675\,000, and the week loses
\$266\,000. Quoting five instruments within one account's budget turns a \$275\,000 week into roughly \$100\,000; a maker adds accounts, venues'
market-maker tiers with larger limits, or fewer instruments, before it adds quotes.

\section{Market-maker programmes and their obligations}

Venues pay makers for liquidity through programmes (One Quant Book 3, chapter 25): a rebate on maker volume if the maker's quotes meet an uptime
requirement, the share of sampled moments at which both sides are within a maximum spread and carry a minimum size. The chapter's maker is two-sided
and within 10 basis points 90\% of the time; its inventory limits, which pull a side when full, are what cost it uptime, and the programme's
requirement is a constraint on its risk management as much as on its technology. Token issuers sign market-making agreements with obligations of
the same kind (One Quant Book 3, chapter 24).

\section{Strategy files}

\begin{strategyfile}{Cross-exchange market making}\label{strat:hf:crypto-high-frequency-trading-on-centralised-venues:xmm}
\sfield{Who pays you, and why} Traders on the less liquid venue, who pay the spread there for immediacy that the deep venue provides more cheaply.
\sfield{Instruments and venues} One coin's perpetual or spot on a quote venue; the deepest spot book as \omterm{def:hf:crypto-high-frequency-trading-on-centralised-venues:hedge}{hedge venue}.
\sfield{Signal} The \omterm{def:hf:crypto-high-frequency-trading-on-centralised-venues:hedge}{hedge venue}'s book as fair value; the quote venue's queue and flow.
\sfield{Sizing and execution} Quote at a half-spread above the hedge's cost; hedge every fill at once; cap inventory to bound the basis risk.
\sfield{Costs} The hedge (1.5 basis points a fill in the model), transfers between venues, custody.
\sfield{How it dies} Venues that cannot be reached with money: Makarov and Schoar's cross-country premia persisted because capital could not move.
\sfield{Horizon, capacity, infrastructure} Milliseconds; accounts and inventory on both venues.
\sfield{Backtest honestly} Both books on one clock; transfer times; the hedge's real fill price.
\sfield{Sources} Makarov and Schoar (2020); this chapter.
\end{strategyfile}

\begin{strategyfile}{Cross-exchange latency taking}\label{strat:hf:crypto-high-frequency-trading-on-centralised-venues:taking}
\sfield{Who pays you, and why} Makers on slower venues whose quotes lag the leading book (chapter 8).
\sfield{Instruments and venues} One coin on several venues; the leader, usually the deepest.
\sfield{Signal} The leader's move against a lagging venue's quotes, beyond fees.
\sfield{Sizing and execution} Take the stale quote; offset on the leader or hold to convergence within inventory limits.
\sfield{Costs} Taker fees on both venues; the risk that the lag is information, not staleness.
\sfield{How it dies} Makers who requote faster: in the model a maker requoting every second leaves nothing; every four seconds, \$135\,000 a week.
\sfield{Horizon, capacity, infrastructure} Milliseconds; co-location where venues offer it.
\sfield{Backtest honestly} The lagging venue's quotes as they were, including the maker's own requote schedule.
\sfield{Sources} Chapter 8; this chapter.
\end{strategyfile}

\begin{strategyfile}{Funding-timestamp trading}\label{strat:hf:crypto-high-frequency-trading-on-centralised-venues:funding}
\sfield{Who pays you, and why} Leveraged traders who hold perpetuals across funding times and pay the rate.
\sfield{Instruments and venues} Perpetuals and their spot hedges; funding at fixed times (00:00, 08:00 and 16:00 UTC on one large venue).
\sfield{Signal} The predicted funding rate from the premium index before each timestamp.
\sfield{Sizing and execution} Hold the side that receives across the timestamp, hedged in spot; unwind after it.
\sfield{Costs} Two legs' fees twice; basis moves around the timestamp as others do the same.
\sfield{How it dies} Crowding: the basis moves before the timestamp and takes the rate with it.
\sfield{Horizon, capacity, infrastructure} Minutes around each timestamp.
\sfield{Backtest honestly} Funding rates as published before the timestamp; the basis in the minutes around it.
\sfield{Sources} One Quant Book 3, chapter 17.
\end{strategyfile}

\begin{strategyfile}{Liquidation-flow absorption}\label{strat:hf:crypto-high-frequency-trading-on-centralised-venues:liquidation}
\sfield{Who pays you, and why} Leveraged traders whose positions the liquidation engine sells regardless of price.
\sfield{Instruments and venues} Perpetuals during cascades; spot on another venue as the hedge.
\sfield{Signal} The basis's depth, the liquidation engine's forced orders, open interest near liquidation prices.
\sfield{Sizing and execution} Absorb once the basis is deep (40 basis points in the model), up to capacity; hedge on spot; hold to the recovery.
\sfield{Costs} Hedging in a falling market; the risk that the dislocation is the new price; the venue's own solvency.
\sfield{How it dies} Cascades deeper than capacity, and venues that auto-deleverage the absorbers' winning positions, as some did on 10 October 2025
when \$19 billion of positions were liquidated in a day.
\sfield{Horizon, capacity, infrastructure} Minutes to hours; balance sheet on several venues.
\sfield{Backtest honestly} Liquidation flows as the engine sent them; auto-deleveraging rules as of the date.
\sfield{Sources} One Quant Book 3, chapter 18; the dated box.
\end{strategyfile}

\begin{strategyfile}{Rate-limit-aware quoting under a programme}\label{strat:hf:crypto-high-frequency-trading-on-centralised-venues:ratelimit}
\sfield{Who pays you, and why} The venue's programme rebate, and clients through the spread.
\sfield{Instruments and venues} Quote venues with request budgets and maker programmes.
\sfield{Signal} The value of a requote: the \omterm{def:hf:fair-value:fair}{fair price}'s move since the last quote against the half-spread.
\sfield{Sizing and execution} Spend requests where the quote is stalest; keep both sides up for uptime; fewer instruments per account.
\sfield{Costs} Stale quotes when the budget binds (\$102\,000 a week at one requote every four seconds in the model).
\sfield{How it dies} Budgets cut by the venue, programmes changed.
\sfield{Horizon, capacity, infrastructure} Seconds; a request governor (\texttt{firm.ratelimit}).
\sfield{Backtest honestly} The venue's limits as of the date (dated box), not the documentation's example values.
\sfield{Sources} One Quant Book 3, chapters 15 and 25.
\end{strategyfile}

\begin{strategyfile}{Spot-perpetual basis market making}\label{strat:hf:crypto-high-frequency-trading-on-centralised-venues:basis}
\sfield{Who pays you, and why} Traders paying for leverage in the perpetual, through funding and the basis.
\sfield{Instruments and venues} A perpetual and its spot, on one or two venues.
\sfield{Signal} The basis against its expected level given the funding rate.
\sfield{Sizing and execution} Quote the basis as one instrument; hedge the legs; carry the position across funding times when the rate pays.
\sfield{Costs} Two legs' fees; basis dislocations in cascades.
\sfield{How it dies} Funding reversals and cascades; the first perpetual, BitMEX's XBTUSD, listed in 2016 with eight-hourly funding, set the template
every venue copied (One Quant Book 3, chapter 17).
\sfield{Horizon, capacity, infrastructure} Hours to weeks.
\sfield{Backtest honestly} Funding and basis histories by venue.
\sfield{Sources} One Quant Book 3, chapter 17.
\end{strategyfile}

\section{Tutorial: around the clock}\label{tut:hf:crypto-high-frequency-trading-on-centralised-venues:tut}

\begin{tutorial}
\textbf{Goal.} Run a week of \omterm{def:hf:crypto-high-frequency-trading-on-centralised-venues:xmm}{cross-exchange market making} with funding, a cascade and a request budget, and decompose the result.
\textbf{End state:} the three figures and the table below.
\begin{tutsteps}
\item \textbf{The quotes} move only when the budget allows; the inventory limits, and absorption in a cascade.
\omcode{code/firm/cryptohft/firm_cryptohft.py}{87}{100}{Requote on the budget's schedule; each side's limit; absorption only when the basis is deep.}\label{lst:hf:crypto-high-frequency-trading-on-centralised-venues:quote}
\item \textbf{The fills and the hedge}: clients, forced sells, arbitrageurs on stale quotes; each fill hedged on spot at once.
\omcode{code/firm/cryptohft/firm_cryptohft.py}{101}{121}{Fills against the quotes, the maker's edge against the perpetual's mid, fees and the hedge's cost.}\label{lst:hf:crypto-high-frequency-trading-on-centralised-venues:fills}
\item \textbf{The week}, the budget sweep and the absorption variants (\texttt{hf\_crypto}); the budget from \texttt{firm.ratelimit}'s governor.
\end{tutsteps}
\textbf{What to change next.} Skew quotes with inventory instead of pulling a side; let the \omterm{def:hf:fair-value:fair}{fair price} come from two \omterm{def:hf:crypto-high-frequency-trading-on-centralised-venues:hedge}{hedge venues}; add transfer
delays and a withdrawal freeze on the quote venue.
\end{tutorial}

\begin{center}
\small
\begin{tabular}{@{}lr@{}}
\toprule
week's component (requote every second, absorbing from 40 bp) & \$\\
\midrule
spread earned on the perpetual & 574\,912\\
lost to stale quotes & 0\\
hedging on spot (half-spread, fee, slippage) & $-464\,310$\\
maker rebates & 143\,728\\
funding & 274\\
cascade (basis recovery) & 20\,842\\
total & 275\,445\\
\bottomrule
\end{tabular}
\end{center}

\section{Build: the crypto cross-exchange maker}\label{bld:hf:crypto-high-frequency-trading-on-centralised-venues:cryptohft}

\begin{build}
\bfield{Purpose} Simulate cross-exchange making with hedging, funding, a liquidation cascade and a request budget.
\bfield{Interface} \texttt{Week(seed, days, sigma\_day, clients\_per\_s, clip, buy\_share, cascade\_day, cascade\_depth\_bp, forced\_per\_s)},
\texttt{run(week, half\_bp, requote\_per\_s, inv\_limit, absorb, absorb\_limit, absorb\_from\_bp, record)} and fee and latency parameters,
\texttt{budget\_from(governor, instruments)}. Built on Book 3's \texttt{firm.perp}, \texttt{firm.ratelimit} and \texttt{firm.mmprogram}.
\bfield{Rules} Dollars; basis points; one-second steps; every fill hedged at once; funding paid by the holder at eight-hour marks.
\bfield{Acceptance tests} \texttt{code/firm/cryptohft/tests/}: the parts add up to the total; fresh quotes are never picked off; stale quotes cost
more and use fewer requests; the budget from the governor's daily cap; absorbing earns more from the cascade than pulling and builds a
multi-million position.
\bfield{Stretch} Inventory skew; two \omterm{def:hf:crypto-high-frequency-trading-on-centralised-venues:hedge}{hedge venues}; transfer delays; auto-deleveraging of the absorber.
\end{build}

\begin{omsources}
\item I.~Makarov, A.~Schoar, Trading and arbitrage in cryptocurrency markets, \emph{Journal of Financial Economics} 135(2), 2020, 293--319.
\item One Quant Book 3, chapters 15, 17, 18 and 25, and their sources for the dated box.
\end{omsources}

\section{Exercises}

\begin{exercise}[$\star$]\label{exo:hf:crypto-high-frequency-trading-on-centralised-venues:1}
A maker quotes 2 basis points either side, receives a 0.5-basis-point maker rebate, and hedges by taking at a 0.5-basis-point half-spread and a
1-basis-point fee. What is its edge per round trip, before slippage?
\end{exercise}

\begin{exercise}[$\star$]\label{exo:hf:crypto-high-frequency-trading-on-centralised-venues:2}
A venue allows 200\,000 orders a day. How many requotes (cancel and new) a second is that for one instrument, and for five?
\end{exercise}

\begin{exercise}[$\star$]\label{exo:hf:crypto-high-frequency-trading-on-centralised-venues:3}
A maker is short \$150\,000 of perpetual at a funding timestamp with a rate of 0.01\%. What does it receive or pay?
\end{exercise}

\begin{exercise}[$\star\star$]\label{exo:hf:crypto-high-frequency-trading-on-centralised-venues:4}
Why does absorbing from the first second of the cascade earn less than waiting until the basis is 40 basis points deep?
\end{exercise}

\begin{exercise}[$\star\star$]\label{exo:hf:crypto-high-frequency-trading-on-centralised-venues:5}
Why does a thin request budget lose money to arbitrageurs even though the maker's spread is unchanged?
\end{exercise}

\begin{exercise}[$\star\star$]\label{exo:hf:crypto-high-frequency-trading-on-centralised-venues:6}
What made Makarov and Schoar's cross-country premia persist?
\end{exercise}

\begin{exercise}[$\star\star\star$]\label{exo:hf:crypto-high-frequency-trading-on-centralised-venues:7}
\emph{Coding.} Rerun the absorption with a 60-basis-point threshold and with a \$2 million limit from the first second. Which does better, and why?
\end{exercise}

\begin{exercise}[$\star\star\star$]\label{exo:hf:crypto-high-frequency-trading-on-centralised-venues:8}
\emph{Find the flaw.} ``Our backtest quotes 20 instruments on one account and requotes each every second; it made \$5 million last month.''
\end{exercise}

\section{Problem: Around the Clock}

\begin{problem}[{Weekend problem --- around the clock}]\label{pb:hf:crypto-high-frequency-trading-on-centralised-venues:1}
A crypto market maker quotes a perpetual on one venue and hedges on another for a week that contains one liquidation cascade.

\medskip\noindent\textbf{Part I --- The structure.}
\begin{enumerate}
\item Define \omterm{def:hf:crypto-high-frequency-trading-on-centralised-venues:xmm}{cross-exchange market making} and the \omterm{def:hf:crypto-high-frequency-trading-on-centralised-venues:hedge}{hedge venue}.
\item What did Makarov and Schoar find?
\item Summarise the dated box.
\item Describe the chapter's week.
\end{enumerate}

\medskip\noindent\textbf{Part II --- Funding and cascades.}
\begin{enumerate}[resume]
\item What is the maker exposed to once hedged?
\item Define \omterm{def:hf:crypto-high-frequency-trading-on-centralised-venues:absorb}{liquidation absorption}.
\item Give the cascade's result for pulling, absorbing at once, from 40 and from 60 basis points.
\item What can go wrong with absorption?
\end{enumerate}

\medskip\noindent\textbf{Part III --- Budgets and programmes.}
\begin{enumerate}[resume]
\item Which of the dated box's limits binds, and at how many requotes a second?
\item Give the week's result at one requote every 1, 2, 4 and 10 seconds.
\item What does a programme's uptime requirement measure, and what costs the maker uptime here?
\item How would you spend a thin budget?
\end{enumerate}

\medskip\noindent\textbf{Part IV --- The verdict.}
\begin{enumerate}[resume]
\item State the \emph{named result}: the cross-exchange maker's P\&L by component over a week with one cascade, and the capture it gives up to stay
inside its request budget.
\item Which component is largest, and which is most fragile?
\item What would a second \omterm{def:hf:crypto-high-frequency-trading-on-centralised-venues:hedge}{hedge venue} add?
\item Which strategy file is most exposed to a venue's failure?
\item How should funding timestamps change the maker's inventory policy?
\item What does the maker owe the venue under a programme?
\item How would you test the model's cascade against the October 2025 record?
\item In one sentence: what does a cross-exchange maker sell?
\end{enumerate}
\end{problem}

\section{Interview questions}

\begin{interviewq}[$\star$ \iqroles{trader}]\label{iq:hf:crypto-high-frequency-trading-on-centralised-venues:1}
Bitcoin is 30 dollars higher on one venue than another. What do you check before trading it?
\end{interviewq}

\begin{interviewq}[$\star\star$ \iqroles{developer}]\label{iq:hf:crypto-high-frequency-trading-on-centralised-venues:2}
Design a request governor for three venues with different limits and a quoting engine that must never be banned.
\end{interviewq}

\begin{interviewq}[$\star\star$ \iqroles{researcher}]\label{iq:hf:crypto-high-frequency-trading-on-centralised-venues:3}
How would you forecast the next funding rate, and how much of it is known before the timestamp?
\end{interviewq}

\begin{interviewq}[$\star\star$ \iqroles{risk}]\label{iq:hf:crypto-high-frequency-trading-on-centralised-venues:4}
Your firm holds \$20 million of coins on a venue that suspends withdrawals. What is at risk, and what should the limits have been?
\end{interviewq}

\begin{interviewq}[$\star\star$ \iqroles{trader}]\label{iq:hf:crypto-high-frequency-trading-on-centralised-venues:5}
Open interest is high and the basis is falling fast. Do you bid the perpetual? How much, and where do you hedge?
\end{interviewq}

\begin{interviewq}[$\star\star\star$ \iqroles{researcher}]\label{iq:hf:crypto-high-frequency-trading-on-centralised-venues:6}
With a requote every $\tau$ seconds, a price volatility $\sigma$ per square-root second and a half-spread $h$, estimate the rate at which an
arbitrageur finds a quote stale by more than $h$.
\end{interviewq}
